\documentclass[10pt]{article}

% ---------- Document Identity (update these only) ----------
\newcommand{\DocStage}{}
\newcommand{\DocTitle}{Your Road to Tier One \textbf{SOF}}
\newcommand{\ProgramName}{182-Week Civilian-to-Selection Readiness Pipeline}
\newcommand{\ProgramSubtitle}{}

% ---------- Page ----------
\usepackage[legalpaper,left=0.5in,right=0.5in,top=0.6in,bottom=0.45in]{geometry}

% ---------- Typography ----------
\usepackage[T1]{fontenc}
\usepackage[utf8]{inputenc}
\usepackage{libertinus}
\usepackage{microtype}
\usepackage{parskip}
\usepackage{ragged2e}
\usepackage{textcomp}
\AtBeginDocument{\color{black}}
\AtBeginDocument{\pagecolor{white}}
\AtBeginDocument{\justifying}

% ---------- Landscape pages ----------
\usepackage{pdflscape}

% ---------- Color ----------
\usepackage[table]{xcolor}
\definecolor{StageBlue}{HTML}{1F4E79}
\definecolor{StageBlueDark}{HTML}{163A5A}
\definecolor{LightGray}{HTML}{F2F4F7}
\definecolor{MidGray}{HTML}{D9DEE5}
\definecolor{RuleGray}{HTML}{C7CED8}
\definecolor{HeaderInk}{HTML}{000000}
\definecolor{Ink}{HTML}{000000}

% ---------- Utility ----------
\usepackage{enumitem}
\sloppy
\setlength{\emergencystretch}{2em}

% ---------- Boxes / UI ----------
\usepackage[most]{tcolorbox}

\tcbset{
  charterTitle/.style={
    enhanced,
    colback=StageBlue!4,
    colframe=StageBlue,
    boxrule=1.25pt,
    arc=3pt,
    left=16pt,right=16pt,top=14pt,bottom=14pt,
    borderline north={3.2pt}{0pt}{StageBlue},
    borderline west={4pt}{0pt}{StageBlue},
    borderline south={1.25pt}{0pt}{StageBlue},
    drop shadow={black!25!white},
  },
  charterRule/.style={
    enhanced,
    colback=LightGray,
    colframe=RuleGray,
    boxrule=0.85pt,
    arc=3pt,
    left=12pt,right=12pt,top=10pt,bottom=10pt,
    borderline west={3pt}{0pt}{StageBlue},
  },
  charterBadge/.style={
    enhanced,
    colback=StageBlue,
    coltext=white,
    colframe=StageBlueDark,
    boxrule=0.6pt,
    arc=2pt,
    left=6pt,right=6pt,top=3pt,bottom=3pt,
  }
}
\newtcbox{\Badge}{nobeforeafter,charterBadge}

% ---------- Lists ----------
\setlist[itemize]{leftmargin=1.5em, itemsep=0.25em, topsep=0.25em}
\setlist[enumerate]{leftmargin=1.8em, itemsep=0.25em, topsep=0.25em}

% ---------- TOC ----------
\usepackage{tocloft}
\setcounter{tocdepth}{3}
\setlength{\cftbeforesecskip}{3pt}
\setlength{\cftbeforesubsecskip}{2pt}
\setlength{\cftbeforesubsubsecskip}{0pt}
\renewcommand{\cftsecfont}{\bfseries}
\renewcommand{\cftsecpagefont}{\bfseries}
\renewcommand{\cftsubsecfont}{\normalfont}
\renewcommand{\cftsubsecpagefont}{\normalfont}
\renewcommand{\cftsubsubsecfont}{\normalfont}
\renewcommand{\cftsubsubsecpagefont}{\normalfont}
\setlength{\cftsecindent}{0pt}
\setlength{\cftsubsecindent}{1.25em}
\setlength{\cftsubsubsecindent}{2.8em}
\setlength{\cftsecnumwidth}{2.1em}
\setlength{\cftsubsecnumwidth}{2.8em}
\setlength{\cftsubsubsecnumwidth}{3.6em}

% Remove the built-in TOC heading so only the custom heading is shown
\makeatletter
\renewcommand{\tableofcontents}{\@starttoc{toc}}
\makeatother

% ---------- Header/Footer: page number only ----------
\usepackage{fancyhdr}
\pagestyle{fancy}
\fancyhf{}
\renewcommand{\headrulewidth}{0pt}
\renewcommand{\footrulewidth}{0pt}
\fancyfoot[C]{\thepage}

% ---------- Section formatting ----------
\usepackage{titlesec}

\titleformat{\section}
  {\Large\bfseries\color{Ink}}
  {\thesection}{0.6em}{}
  [\vspace{0.25em}\textcolor{StageBlue}{\rule{\linewidth}{0.8pt}}]

\titleformat{\subsection}
  {\large\bfseries\color{Ink}}
  {\thesubsection}{0.6em}{}

\titleformat{\subsubsection}
  {\normalsize\bfseries\color{Ink}}
  {\thesubsubsection}{0.6em}{}

\titlespacing*{\section}{0pt}{0.95em}{0.35em}
\titlespacing*{\subsection}{0pt}{0.65em}{0.25em}
\titlespacing*{\subsubsection}{0pt}{0.55em}{0.2em}

% ---------- Table engine ----------
\usepackage{tabularray}
\UseTblrLibrary{booktabs}

% ---------- tabularray: GLOBAL table treatment ----------
\SetTblrInner{
  cells = {fg=black, bg=white, valign=t},
}

% ---------- Header helpers ----------
\newcommand{\Hdr}[1]{\shortstack[c]{\strut #1\strut}}

% ---------- Lines ----------
\newcommand{\TitleLine}{\textcolor{StageBlue}{\rule{0.98\linewidth}{0.95pt}}}
\newcommand{\SubTitleLine}{\textcolor{RuleGray}{\rule{0.98\linewidth}{0.65pt}}}

% ---------- Continued on next page ----------
\newcommand{\ContinuedOnNextPageInline}{%
  \par
  \vfill
  \noindent\textcolor{RuleGray}{\rule{\linewidth}{1pt}}\vspace{-2pt}
  \noindent\hfill\textit{Continued on next page}\par\vspace{-10pt}
  \noindent\textcolor{RuleGray}{\rule{\linewidth}{1pt}}
  \clearpage
}

% ---------- PDF hyperlinks ----------
\usepackage[hidelinks]{hyperref}
\hypersetup{
  colorlinks=false,
  pdfborder={0 0 0},
  linktoc=all,
  pdftitle={\DocTitle},
  pdfauthor={\ProgramName}
}

% =========================================================
\begin{document}

\section*{10.C.1 Core Doctrine — Program-Wide Rules}
\phantomsection
\addcontentsline{toc}{section}{10.C.1 Core Doctrine — Program-Wide Rules}

\begin{itemize}
\item Baseline plus session top-up model:
  \begin{itemize}
  \item meet a conservative daily baseline intake band, then add session-specific fluids based on duration and environment.
  \end{itemize}
\item Sip model only:
  \begin{itemize}
  \item sip regularly; no forced chugging; no “catch-up dumping” large volumes in short time windows.
  \end{itemize}
\item Avoid extremes:
  \begin{itemize}
  \item dehydration practices and water cuts are prohibited,
  \item overhydration is a safety risk; do not drink past comfort cues.
  \end{itemize}
\item Electrolytes are condition-dependent:
  \begin{itemize}
  \item electrolytes are not mandatory daily intake; they are activated only when trigger conditions are met.
  \end{itemize}
\item No novelty near gates and test weeks:
  \begin{itemize}
  \item no new electrolyte products, concentrations, or timing practices inside the final 7–10 days before major gate tests,
  \item no caffeine pattern changes within 14 days of gate windows (caffeine posture is governed elsewhere; electrolytes must not be used as stimulant substitutes).
  \end{itemize}
\item If a change is unavoidable due to logistics:
  \begin{itemize}
  \item log it clearly (date, reason, what changed),
  \item interpret outcomes conservatively (do not claim comparable performance evidence if conditions/inputs drift).
  \end{itemize}
\end{itemize}

\newpage

\section*{10.C.2 Baseline Daily Hydration Posture}
\phantomsection
\addcontentsline{toc}{section}{10.C.2 Baseline Daily Hydration Posture}

Baseline intake framework (locked):

\begin{itemize}
\item Minimum floor: 30 mL/kg/day.
\item Typical range: 35–45 mL/kg/day.
\item Present as ranges; variability is expected.
\end{itemize}

Hybrid bodyweight reference rule (locked):

\begin{itemize}
\item Phase 00–Phase 01 baseline calculations use current bodyweight reference: 266 lb (121 kg).
\item Transition toward target bodyweight reference of 180 lb (82 kg) by Phase 05 for baseline calculations (do not apply abrupt step changes; transition is gradual and logged).
\end{itemize}

Conversions for current bodyweight example (non-clinical estimate; Phase 00–01, 121 kg):

\begin{itemize}
\item Floor (30 mL/kg/day): 121 × 30 = 3630 mL/day = 3.63 L/day = 123 oz/day.
\item Typical band (35–45 mL/kg/day):
  \begin{itemize}
  \item 121 × 35 = 4235 mL/day = 4.24 L/day = 143 oz/day.
  \item 121 × 45 = 5445 mL/day = 5.45 L/day = 184 oz/day.
  \end{itemize}
\end{itemize}

Education-only heuristics (non-diagnostic):
\begin{itemize}
\item thirst and urine color are rough awareness cues only; they do not override symptom-based stop posture or safety decisions.
\end{itemize}

Distribution rule (binding):
\begin{itemize}
\item spread intake across the day,
\item avoid large late-day boluses,
\item avoid rapid catch-up drinking.
\end{itemize}

\begingroup
\scriptsize
\setlength{\tabcolsep}{2pt}
\renewcommand{\arraystretch}{1.12}

\begin{longtblr}{
  width=\linewidth,
  colspec = {
    X[l,2.55]
    X[l,2.65]
    X[l,4.80]
  },
  rowhead=1,
  hlines,
  vlines,
  cells = {hsep=2pt, vsep=6pt, halign=l, valign=t},
  row{odd} = {bg=white},
  row{even} = {bg=LightGray},
  row{1} = {bg=StageBlue!15, fg=black, font=\bfseries, halign=l, valign=m},
}
\Hdr{Condition} &
\Hdr{Baseline Intake Posture} &
\Hdr{Notes} \\
Normal indoor day, low sweat &
30–35 mL/kg/day floor-to-low band &
Use distributed sips; avoid late-day boluses. \\
Typical training day, moderate sweat &
35–45 mL/kg/day typical band &
Baseline is not a session replacement; session top-ups still apply. \\
Hot day / heavy sweat day &
45 mL/kg/day band bias plus session top-ups &
Do not force drink; use sip model; activate electrolyte posture when triggers met. \\
Cold day (muted thirst) &
Maintain baseline band; schedule sips &
Cold suppresses thirst; underdrinking risk increases; warm fluids allowed if feasible. \\
GI upset day (mild) &
Maintain baseline cautiously; prioritize small sips &
If severe GI distress or vomiting/diarrhea persists, stop and seek evaluation posture. \\
\end{longtblr}

\endgroup

\newpage

\section*{10.C.3 Session Hydration Rules — By Duration and Environment}
\phantomsection
\addcontentsline{toc}{section}{10.C.3 Session Hydration Rules — By Duration and Environment}

Session types and per-hour fluid bands (locked):

\begin{itemize}
\item Cool or indoor: 8–16 oz/hr.
\item Hot or outdoor: 16–24 oz/hr.
\item Sip model only; do not force drinking.
\item Hard ceiling (locked): do not exceed 32 oz/hr unless clinically directed.
\end{itemize}

Define pre/during/post posture by session type:

\begin{itemize}
\item $\leq$60 minutes, low to moderate effort:
  \begin{itemize}
  \item Pre: baseline hydration maintained; small pre-session sips if needed.
  \item During: minimal sips as desired; band still applies but typically low end.
  \item Post: resume baseline; avoid rapid catch-up.
  \end{itemize}
\item 60–90 minutes:
  \begin{itemize}
  \item Pre: ensure baseline plus modest top-up before start (avoid bolus).
  \item During: follow per-hour band by environment; sip model.
  \item Post: resume baseline; monitor symptoms.
  \end{itemize}
\item $\geq$90 minutes or prolonged outdoor / ruck exposure:
  \begin{itemize}
  \item Pre: baseline + modest top-up; confirm access plan (bottle/reservoir capacity per 7.13).
  \item During: use per-hour band; activate electrolyte posture when triggers met.
  \item Post: resume baseline; do not chug.
  \end{itemize}
\item Hot/humid:
  \begin{itemize}
  \item bias to higher end of per-hour band; prioritize safety; reslot tests if heat index triggers are exceeded.
  \end{itemize}
\item Cold:
  \begin{itemize}
  \item schedule sips; cold-sweat risk exists; maintain baseline.
  \end{itemize}
\end{itemize}

\begingroup
\scriptsize
\setlength{\tabcolsep}{2pt}
\renewcommand{\arraystretch}{1.12}

\begin{longtblr}{
  width=\linewidth,
  colspec = {
    X[l,1.90]
    X[l,2.35]
    X[l,2.55]
    X[l,2.35]
    X[l,2.00]
    X[l,2.80]
  },
  rowhead=1,
  hlines,
  vlines,
  cells = {hsep=2pt, vsep=6pt, halign=l, valign=t},
  row{odd} = {bg=white},
  row{even} = {bg=LightGray},
  row{1} = {bg=StageBlue!15, fg=black, font=\bfseries, halign=l, valign=m},
}
\Hdr{Session Type} &
\Hdr{Pre-Session} &
\Hdr{During} &
\Hdr{Post-Session} &
\Hdr{Electrolyte Posture} &
\Hdr{Notes and Safety} \\
$\leq$60 min, low–moderate &
Maintain baseline; optional 4–8 oz sips if thirsty &
Sip as desired; typically within 8–16 oz/hr cool/indoor &
Resume baseline; avoid catch-up chugging &
Typically none &
Do not force drinking; stop for dizziness, nausea, unusual distress. \\
60–90 min &
Baseline + modest top-up (e.g., small sips over 30–60 min prior) &
8–16 oz/hr cool/indoor; 16–24 oz/hr hot/outdoor &
Resume baseline; monitor cramps/headache &
Consider if heavy sweat or heat exposure &
Avoid novelty products near gates; log environment and intake estimate. \\
$\geq$90 min or prolonged outdoor / ruck &
Baseline + planned access (bottle/reservoir capacity) &
8–16 oz/hr cool/indoor; 16–24 oz/hr hot/outdoor; do not exceed 32 oz/hr &
Resume baseline; spread intake; include food as normal (no new protocol) &
Recommended when triggers met ($\geq$90 min, heavy sweat, heat, ruck) &
Overhydration risk rises with prolonged sessions; electrolyte posture reduces “water-only overload” risk when appropriately triggered. \\
Hot/humid exposure &
Baseline plus earlier pre-session sips &
Bias 16–24 oz/hr; sip model; ceiling 32 oz/hr &
Resume baseline; cool down; monitor symptoms &
Recommended when triggered &
Heat triggers can invalidate testing; reslot per Stage 3 when unsafe. \\
Cold exposure (muted thirst) &
Maintain baseline; warm fluids allowed &
Schedule sips; keep within 8–16 oz/hr unless heavy clothing sweat &
Resume baseline; avoid late-day bolus &
Consider if prolonged outdoor with sweat &
Cold suppresses thirst; underdrinking risk increases; do not ignore symptoms. \\
\end{longtblr}

\endgroup

\section*{10.C.4 Electrolyte Operating Rules — Conservative, Non-Prescriptive}
\phantomsection
\addcontentsline{toc}{section}{10.C.4 Electrolyte Operating Rules — Conservative, Non-Prescriptive}

Activation threshold (locked): electrolytes are recommended only if one or more apply:

\begin{itemize}
\item session duration $\geq$90 minutes,
\item heavy sweat or heat exposure,
\item ruck or prolonged outdoor exposure.
\end{itemize}

Sodium guidance bands (locked; no product dosing instructions):

\begin{itemize}
\item 300–600 mg sodium per hour in moderate conditions.
\item 600–1,000 mg sodium per hour in hot conditions or heavy sweat.
\end{itemize}

Potassium and magnesium bands (locked; conservative; aligned to sodium timing):

\begin{itemize}
\item Potassium: 100–200 mg per hour (moderate); 200–300 mg per hour (hot/heavy sweat).
\item Magnesium: 25–50 mg per hour (moderate); 50–75 mg per hour (hot/heavy sweat).
\end{itemize}

Safety cautions (binding):
\begin{itemize}
\item Do not megadose.
\item If known or suspected kidney disease, clinician-directed fluid/sodium restrictions, heart failure, uncontrolled hypertension, or similar clinician constraints exist: do not self-direct electrolyte strategy; clinician oversight is required.
\end{itemize}

Product selection posture (generic, non-brand):

\begin{itemize}
\item transparent labeling of sodium/potassium/magnesium amounts,
\item no proprietary blends masking amounts,
\item avoid stimulant-combined electrolyte products; electrolytes are not stimulant substitutes; cross-reference Stage 7.05 boundaries by \textbf{ID} only.
\end{itemize}

Mixing posture (binding):

\begin{itemize}
\item follow label mixing instructions,
\item do not concentrate beyond label guidance,
\item do not create “double-strength” mixtures in protected windows.
\end{itemize}

\begingroup
\scriptsize
\setlength{\tabcolsep}{2pt}
\renewcommand{\arraystretch}{1.12}

\begin{longtblr}{
  width=\linewidth,
  colspec = {
    X[l,2.10]
    X[l,3.45]
    X[l,3.35]
    X[l,3.55]
  },
  rowhead=1,
  hlines,
  vlines,
  cells = {hsep=2pt, vsep=6pt, halign=l, valign=t},
  row{odd} = {bg=white},
  row{even} = {bg=LightGray},
  row{1} = {bg=StageBlue!15, fg=black, font=\bfseries, halign=l, valign=m},
}
\Hdr{Trigger} &
\Hdr{Electrolyte Response} &
\Hdr{What to Log} &
\Hdr{When to Stop or Escalate} \\
$\geq$90 min session &
Apply moderate band: Na 300–600 mg/hr; K 100–200 mg/hr; Mg 25–50 mg/hr &
Session duration, environment, electrolyte Y/N, estimated bands used &
Severe GI distress, dizziness, confusion, chest pressure, near-syncope → stop and seek evaluation posture \\
Hot/heavy sweat &
Apply hot band: Na 600–1,000 mg/hr; K 200–300 mg/hr; Mg 50–75 mg/hr &
Heat/humidity notes; sweat indicators; intake estimate &
Heat illness signs, severe headache, vomiting, weakness → stop and seek evaluation posture; reslot tests \\
Ruck/prolonged outdoor &
Treat as trigger; apply moderate or hot band based on conditions &
Dry load context noted separately in training logs; intake estimate &
Gait-altering pain plus systemic symptoms → stop exposure and reslot; seek evaluation if red flags \\
No trigger present &
No electrolytes by default &
Baseline fluids only; electrolytes N &
If cramping persists with normal hydration posture, downshift and review; do not escalate to megadoses \\
\end{longtblr}

\endgroup

\newpage

\section*{10.C.5 Heat and Cold Safety Rules — Hydration as Risk Control}
\phantomsection
\addcontentsline{toc}{section}{10.C.5 Heat and Cold Safety Rules — Hydration as Risk Control}

Numeric decision triggers (locked; conservative) and required actions; safety overrides schedule and validity.

Heat / Heat Index triggers:

\begin{itemize}
\item Heat Index 85–89°F:
  \begin{itemize}
  \item caution posture; shorten exposure; indoor option ready; increase monitoring and sips.
  \end{itemize}
\item Heat Index 90–94°F:
  \begin{itemize}
  \item indoor substitution strongly preferred for longer sessions and any near-max efforts.
  \end{itemize}
\item Heat Index 95–99°F:
  \begin{itemize}
  \item reslot gate tests; indoor substitution for training; avoid prolonged outdoor exposure.
  \end{itemize}
\item Heat Index 100°F or higher:
  \begin{itemize}
  \item do not test; do not execute prolonged outdoor exposure; safety overrides schedule.
  \end{itemize}
\end{itemize}

Cold triggers:

\begin{itemize}
\item Air temperature 10–19°F OR windchill 0–9°F:
  \begin{itemize}
  \item hydration remains required; schedule sips; recognize cold-sweat risk under layers.
  \end{itemize}
\item Windchill below 0°F:
  \begin{itemize}
  \item indoor substitution preferred for long sessions and any test efforts; document substitution.
  \end{itemize}
\end{itemize}

Visibility/storm triggers overriding outdoor plans:

\begin{itemize}
\item lightning within area,
\item active ice glaze,
\item whiteout or near-zero visibility,
\item official air quality warning,
\item → substitute indoors or reslot per Stage 3 adjustment doctrine.
\end{itemize}

\begingroup
\scriptsize
\setlength{\tabcolsep}{2pt}
\renewcommand{\arraystretch}{1.12}

\begin{longtblr}{
  width=\linewidth,
  colspec = {
    X[l,2.15]
    X[l,2.675]
    X[l,2.95]
    X[l,2.85]
    X[l,3.30]
  },
  rowhead=1,
  hlines,
  vlines,
  cells = {hsep=2pt, vsep=6pt, halign=l, valign=t},
  row{odd} = {bg=white},
  row{even} = {bg=LightGray},
  row{1} = {bg=StageBlue!15, fg=black, font=\bfseries, halign=l, valign=m},
}
\Hdr{Situation} &
\Hdr{Numeric Trigger Example} &
\Hdr{Immediate Action} &
\Hdr{Next-Step Adjustment} &
\Hdr{Notes} \\
Moderate heat caution &
Heat Index 85–89°F &
Proceed only with caution; shorten exposure; plan indoor fallback &
Use sip model; consider electrolytes if duration $\geq$90 min &
Log heat category; do not force drinking; monitor symptoms \\
High heat &
Heat Index 90–94°F &
Prefer indoor substitution for long sessions/near-max efforts &
Reslot any gate-grade testing &
Heat can invalidate validity and increase risk \\
Very high heat &
Heat Index 95–99°F &
Reslot gate tests; substitute indoors &
Maintain hydration baseline; avoid prolonged outdoor &
Safety and validity override schedule \\
Extreme heat &
Heat Index $\geq$100°F &
Do not test; do not do prolonged outdoor exposure &
Indoor-only or reslot &
Stop posture escalates on symptoms \\
Cold exposure &
Air 10–19°F or windchill 0–9°F &
Proceed only with scheduled sips &
Maintain baseline; avoid dehydration by muted thirst &
Cold does not remove hydration requirement \\
Severe cold &
Windchill <0°F &
Indoor substitution preferred for long sessions/tests &
Reslot test efforts &
Document substitution; protect safety \\
Visibility/storm &
Lightning/ice glaze/whiteout/air quality warning &
Stop outdoor plan immediately &
Substitute indoors or reslot per Stage 3 &
Do not attempt to “salvage” a test under unsafe conditions \\
\end{longtblr}

\endgroup

\newpage

\section*{10.C.6 Test Week Hydration Posture — Validity Protection}
\phantomsection
\addcontentsline{toc}{section}{10.C.6 Test Week Hydration Posture — Validity Protection}

Binding posture:

\begin{itemize}
\item Keep hydration routine stable during test weeks and the final 7–10 days before major gate tests.
\item Avoid new electrolyte products or changes in concentration or timing practices in protected windows.
\item Ensure adequate fluids in the 24 hours prior using baseline plus conservative session top-ups; avoid large late-day boluses.
\item Log environment conditions affecting comparability and safety.
\item Reslot tests when conditions invalidate safety or validity per Stage 3 adjustment doctrine.
\end{itemize}

\begingroup
\scriptsize
\setlength{\tabcolsep}{2pt}
\renewcommand{\arraystretch}{1.12}

\begin{longtblr}{
  width=\linewidth,
  colspec = {
    X[l,1.70]
    X[l,2.55]
    X[l,3.60]
    X[l,3.35]
    X[l,3.40]
  },
  rowhead=1,
  hlines,
  vlines,
  cells = {hsep=2pt, vsep=6pt, halign=l, valign=t},
  row{odd} = {bg=white},
  row{even} = {bg=LightGray},
  row{1} = {bg=StageBlue!15, fg=black, font=\bfseries, halign=l, valign=m},
}
\Hdr{Week Type} &
\Hdr{Goal} &
\Hdr{Allowed Adjustments} &
\Hdr{Not Allowed} &
\Hdr{Notes} \\
Deload/Test week &
Preserve validity and reduce fatigue contamination &
Calories set to maintenance per Stage 10.B; hydration baseline stable; session top-ups as needed &
New electrolyte product; new concentration; new timing practice; new caffeine pattern &
Treat electrolytes as controlled variables; stability is prioritized over “optimization” \\
Test day(s) &
Stable output and safety &
If already using electrolytes, use the same product type and same mixing/timing; sip model &
First-time electrolyte use; double-strength mixes; chugging &
Log conditions and intake estimate; reslot if heat/cold triggers invalidate safety \\
7–10 days pre-gate &
Controlled variable stability &
Hold routine; minor logistical adjustments only if unavoidable and logged &
Any structural change to hydration/electrolyte routine &
Interpret outcomes conservatively if any unavoidable change occurs \\
\end{longtblr}

\endgroup

\newpage

\section*{10.C.7 Logging Requirements — Hydration Fields}
\phantomsection
\addcontentsline{toc}{section}{10.C.7 Logging Requirements — Hydration Fields}

Fields to record aligned to Stage 2.E posture (operational; no schema rewrite):

\begin{itemize}
\item Daily hydration estimate:
  \begin{itemize}
  \item approximate oz/day or L/day OR low/medium/high category, recorded consistently.
  \end{itemize}
\item Electrolyte use:
  \begin{itemize}
  \item Y/N plus generic product-type note (no brands) and “moderate band” vs “hot band” intent.
  \end{itemize}
\item Session hydration estimate:
  \begin{itemize}
  \item approximate pre/during/post volume or rough total.
  \end{itemize}
\item Environment notes relevant to hydration:
  \begin{itemize}
  \item hot/humid/cold; indoor/outdoor; unusual sweat indicators.
  \end{itemize}
\item Symptom flags:
  \begin{itemize}
  \item cramping; dizziness; headache; unusual fatigue; nausea or GI distress.
  \end{itemize}
\end{itemize}

\begingroup
\scriptsize
\setlength{\tabcolsep}{2pt}
\renewcommand{\arraystretch}{1.12}

\begin{longtblr}{
  width=\linewidth,
  colspec = {
    Q[l,wd=0.5in]
    Q[l,wd=.5in]
    Q[r,wd=0.5in]
    X[l,2.10]
    X[l,1.55]
    Q[l,wd=1.00in]
    X[l,1.20]
    X[l,1.50]
  },
  rowhead=1,
  hlines,
  vlines,
  cells = {hsep=2pt, vsep=6pt, halign=l, valign=t},
  row{odd} = {bg=white},
  row{even} = {bg=LightGray},
  row{1} = {bg=StageBlue!15, fg=black, font=\bfseries, halign=l, valign=m},
}
\Hdr{Date} &
\Hdr{Phase/Week} &
\Hdr{Session \#} &
\Hdr{Environment} &
\Hdr{Fluid Estimate} &
\Hdr{Electrolytes} &
\Hdr{Symptoms} &
\Hdr{Notes} \\
&
&
&
&
&
&
&
\\
\end{longtblr}

\endgroup

\newpage

\section*{10.C.8 Common Failure Modes and Controls — Hydration-Specific}
\phantomsection
\addcontentsline{toc}{section}{10.C.8 Common Failure Modes and Controls — Hydration-Specific}

\begin{itemize}
\item Underdrinking due to access limits → pre-plan bottle/reservoir capacity and refill points via 7.13; do not start sessions dry.
\item Underdrinking in cold because thirst is muted → schedule sips; warm fluids if feasible; maintain baseline.
\item Overdrinking to catch up → sip model; spread intake across day; avoid rapid chugging.
\item Overdrinking plain water during long/hot sessions → activate conservative electrolyte posture when triggered; avoid forced water-only intake.
\item Treating electrolytes as mandatory daily intake → electrolytes are condition-dependent; use only when triggers apply.
\item Introducing new electrolyte products or concentrations on test day → prohibited; no novelty near gates/test weeks.
\item Ignoring early warning signs → downshift, move indoors, review baseline plus top-up; log signals.
\item Poor logging (missing environment/symptoms) → hydration fields are mandatory in session notes; missing fields reduce interpretability and safety controls.
\item Failure to reslot tests in unsafe heat/illness conditions → apply Stage 3 reslot posture; safety and validity override schedule.
\item Using stimulant-combined electrolytes as an energy substitute → prohibited; electrolytes are not stimulant substitutes; respect 7.05 boundaries by \textbf{ID} only.
\end{itemize}

\newpage

\section*{10.C.9 Escalation Triggers — Non-Medical}
\phantomsection
\addcontentsline{toc}{section}{10.C.9 Escalation Triggers — Non-Medical}

Stop-and-seek-evaluation triggers (non-diagnostic safety posture):

\begin{itemize}
\item confusion or altered mental status,
\item fainting or near-fainting,
\item chest pain or pressure,
\item severe dizziness or inability to stand/walk normally,
\item severe headache especially with heat exposure,
\item unusual shortness of breath not explained by exertion level,
\item persistent vomiting or diarrhea,
\item signs consistent with suspected heat illness especially when combined with the above.
\end{itemize}

Explicit statement: this list is not a diagnosis; it is a safety escalation posture.

\newpage

\section*{10.C.10 Operating System Integrity Checks}
\phantomsection
\addcontentsline{toc}{section}{10.C.10 Operating System Integrity Checks}

\begin{itemize}
\item Baseline hydration posture exists as ranges with a minimum floor (30 mL/kg/day) and typical range (35–45 mL/kg/day) using hybrid bodyweight reference rules across phases.
\item Baseline plus session top-up model is specified and usable without specialized tools, using 7.13 container capability.
\item Session rules exist for $\leq$ 60, 60–90, and $\geq$90 minutes, with hot/humid and cold modifiers and locked per-hour bands.
\item Electrolytes are condition-dependent; activation thresholds are explicit ($\geq$90 minutes, heavy sweat/heat, ruck/prolonged outdoor).
\item Sodium numeric bands exist (300–600 mg/hr; 600–1,000 mg/hr) and are conservative and non-prescriptive.
\item Potassium and magnesium numeric bands exist with explicit “no megadose” and clinician-caution posture for kidney disease/clinician restrictions.
\item Overhydration ceiling exists: do not exceed 32 oz/hr unless clinically directed; do not chug posture is explicit.
\item Numeric heat and cold triggers exist and explicitly reslot tests per Stage 3 when conditions invalidate safety or validity.
\item Test-week stability posture exists; no novelty near gates/test weeks is enforced for hydration/electrolyte routines.
\item Logging fields align to Stage 2.E conventions and are implementable under start-from-zero posture, including environment notes and symptom flags.
\item Any upstream conflict is labeled "Requires Stage 8 review" with no silent edits.
\end{itemize}

\end{document}